\section{Gaussian tensor networks}

While gaussian chaoses are defined by tensor products, let us define general Gaussian tensor networks on hypergraphs $\graph$ by tensors
\begin{align*}
    \gausscoreofat{\graph}{\nodevariables}
\end{align*}
with coordinates
\begin{align*}
    \gausscoreofat{\graph}{\indexedcatvariableof{\nodes}}
    = \prod_{\edgein}\gausscoreofat{\edge}{\indexedcatvariableof{\edge}} \, .
\end{align*}

Especially we are interested on moment bounds on
\begin{align*}
    \contraction{\hypercorewith,\gausscoreofat{\graph}{\nodevariables}}
\end{align*}
where $[\catorder]\subset\nodes$.
Intuitively, averaging over hidden variables improves the concentration behavior of Gaussian tensor networks, due to the law of large numbers.
Here we need to show effects of averaging over multiple hidden variables.

\subsection{Reduction to Gaussian chaoses}

The trick we apply is that any contracted Gaussian tensor network can be related to a Gaussian chaos on a mutilated hypergraph.
This mutilation involves introduction of additional edges decorated by delta tensors.

To be more precise we define a new hypergraph
\begin{align*}
    \nodes \cup \bigcup_{\edge\in\edges}\bigcup_{\node\in\edge} \{(\edge,\node)\}
\end{align*}
and edges
\begin{align*}
    \{\{(\edge,\node) \wcols \node\in\edge\} \wcols \edgein \} % for the original edges
    \cup \{\{(\edge,\node) \wcols \node\in\edge\}\cup\{\node\} \wcols \nodein \} % for the delta tensors
\end{align*}
where the first set of edges is decorated as before and the second by delta tensors.
Then
\begin{align*}
    \contraction{\hypercorewith,\gausscoreofat{\graph}{\nodevariables}}
    = \contraction{\{\hypercorewith\}\cup\{\delta[(\catvariableof{\edge,\node}:\node\in\edge),\catvariableof{\node}]\}
        \cup \{\gausscoreofat{\edge}{\catvariableof{\edge,\edge}}\wcols\edgein\}
    }
\end{align*}

We notice that this is a Gaussian chaos to the structure tensor
\begin{align*}
    \contractionof{\{\hypercorewith\}\cup\{\delta[(\catvariableof{(\edge,\node)}:\node\in\edge),\catvariableof{\node}]\}}{
        \catvariableof{(\node,\edge)\wcols \edgein,\node\in\edge}
    } \, .
\end{align*}

To show moment bounds, we are left to calculate the $J$-injective norms of this structure tensor.

\subsection{Properties of $J$-injective norms}

\begin{itemize}
    \item Tensor product property: Multiplicativity of $J$-injective norms into those to factors.
    \item Delta tensor property (more general: ODECO).
\end{itemize}